% Generated draft: requires mathematical and source review.
\documentclass[11pt]{article}
\usepackage[margin=25mm]{geometry}
\usepackage{fontspec}
\setmainfont{Latin Modern Roman}
\usepackage{amsmath,amssymb,mathtools,unicode-math}
\setmathfont{Latin Modern Math}
\usepackage{enumitem,needspace,xurl,hyperref}
\setlength{\parindent}{0pt}
\setlength{\parskip}{5pt}
\setlength{\emergencystretch}{3em}
\newcommand{\N}{\mathbb{N}}
\newcommand{\Z}{\mathbb{Z}}
\newcommand{\Q}{\mathbb{Q}}
\newcommand{\R}{\mathbb{R}}
\newcommand{\C}{\mathbb{C}}
\newcommand{\F}{\mathbb{F}}
\newcommand{\A}{\mathbb{A}}
\newcommand{\PP}{\mathbb{P}}
\newcommand{\E}{\mathbb{E}}
\newcommand{\eps}{\varepsilon}
\newcommand{\dd}{\,\mathrm{d}}
\begin{document}
\section*{UnsolvedMath ID: 10600002}
\textbf{Source problem code:} catalogue entry
\subsection{Definitions and mathematical statement}
In the usual diagrammatic model, a flat virtual knot is represented by an immersed circle with flat crossings (having no over/under information) and virtual crossings, modulo flat Reidemeister moves and virtual detour moves. The source also calls these objects virtual strings; its orientation conventions are not specified here. A combinatorial invariant is a function on such diagrams that is unchanged by the allowed moves, and therefore descends to their equivalence classes.

Problem-list item 2 asks for new combinatorial invariants of flat virtual knots and for a better understanding of the flat biquandle algebra, including constructions that detect flat linking information beyond linking numbers modulo $2$. These are requests for new invariants and algebraic understanding, rather than a single specified classification theorem.

The source asserts that the flat biquandle algebra is isomorphic to a Weyl algebra and has no nontrivial finite-dimensional representations. Its presentation, coefficient ring, characteristic, and representation conventions are absent from the supplied excerpt, so this assertion must retain that qualification. For comparison only, the standard first Weyl algebra over a field $k$ is $A_1(k)=k\langle x,y\rangle/(xy-yx-1)$; the familiar obstruction to nonzero finite-dimensional unital representations assumes characteristic zero. The supplied text does not establish that this particular presentation and these hypotheses are its intended ones. It cites a simple affine biquandle used for the Affine Index polynomial as going beyond the usual restrictions, but supplies neither its operations nor an exact theorem relating it to the asserted obstruction. The problem therefore does not ask for a representation contradicting an established nonexistence statement under the same hypotheses.
\subsection{Short English statement}
Construct new invariants of flat virtual knots and clarify the flat biquandle algebra and its useful realizations, especially for linking information beyond parity. The source's Weyl-algebra and representation claims require missing coefficient and presentation conventions.
\subsection{Sources}
\begin{itemize}
\item Local source registry: \url{data/unsolvedmath/problems.json}, permanent UnsolvedMath ID $10600002$, exact source code \texttt{catalogue entry}, ``Virtual-knot problem 2 --- Flat Virtuals.''
\item \emph{Unsolved Problems in Virtual Knot Theory and Combinatorial Knot Theory}, revised survey, dedicated problem-list item 2: \url{https://arxiv.org/abs/1409.2823}. Supplied access date: 2026-07-29; extraction: source-tex. The source list contains 60 top-level items, with compound items retained intact.
\item The registry mentions semiquandle colorings, $u$-polynomials, polynomial flat-link invariants, flat virtual braid representations, and works titled ``Connected sum and crossing numbers of flat virtual knots'' and ``Equivalence of flat-virtual diagrams.'' Its arXiv strings have placeholder endings \texttt{xxxx}; these are incomplete references and cannot serve as verified citations. No precise reference for the Affine Index polynomial construction is supplied in this record.
\end{itemize}
\end{document}
